\section{数组、双指针和滑动窗口章正文案例推导补强}

这一节补齐本章正文出现过的 LeetCode 案例推导。阅读顺序统一按“题目说明、样例入口、暴力基线、暴力过程、重复工作、优化条件、相邻状态、状态复用、指针和字段、代码落点、正确性、边界实验、复杂度变化、迁移追问”展开；目的不是新增题量，而是把正文中的案例都讲成可复述、可证明、可迁移的解题链。

\subsection{LeetCode 53 最大子数组和}

\textbf{题目说明。} LeetCode 53 最大子数组和给定整数数组 nums，要求找一个非空连续子数组，使元素和最大，并返回这个最大和。

\textbf{样例入口。} 样例 nums=[-2, 1, -3, 4, -1, 2, 1, -5, 4]，输出 6；连续段 [4, -1, 2, 1] 的和为 6。

\textbf{暴力基线。} 暴力枚举所有连续子数组，累计每段和并更新最大值；优化一点可以固定左端、右端扩展时维护当前和。

\textbf{暴力过程。} 以 [-2, 1, -3, 4, -1, 2, 1] 为例，暴力会重复计算很多以 4 开头或以 2 结尾的子段。真正关键是当前数要不要接上前一个后缀。

\textbf{重复工作。} 题面模型：给定整数数组，返回一个非空连续子数组能够取得的最大元素和。暴力版的慢点要落到本题：浪费在重复求每个后缀和。对于以 i 结尾的最优连续段，历史信息只需要一个值：以 i-1 结尾的最大后缀和。后面的优化只消掉这类重复工作，不能改变答案集合。

\textbf{优化条件。} 本题的关键观察是：用 [-2, 3] 和 [2, 3] 对比，推出当前位置只在单独开始与接上旧后缀之间取较大值。这句话必须能翻译成可维护状态；如果题目条件改变后观察不再成立，就不能继续套原来的写法。

\textbf{相邻状态。} 规律是当前位置只需要比较 nums[i] 与 bestEndingHere + nums[i]，全局答案再和 bestEndingHere 比较；全负数组要求答案初始化为第一个数，不能初始化为 0。把这个特例继续拆开看：先问暴力搜索里哪些不同路径到达同一个后续问题，再给这个后续问题命名为状态。状态必须包含未来决策需要的全部信息，转移只从更小且已经算清的状态来。

\textbf{状态复用。} bestEndingHere 表示以当前位置结尾的最大子数组和；answer 表示全局最大值。当前位置要么单独开始，要么接上旧后缀。手算时只改被新输入影响的那一小部分，而不是回头重算完整候选。

\textbf{指针和字段。} 状态字段要能说清：current=bestEndingHere，answer 是全局最优；nums[i] 与 current+nums[i] 比较决定是否丢弃旧后缀。手算时按这个协议跟踪：遇到 4 之前，旧后缀为负，接上反而变差，所以从 4 重新开始；之后 -1、2、1 继续接上得到 6。

\textbf{代码落点。} 关键代码是 current=max(nums[i], current+nums[i]); answer=max(answer, current)。answer 必须初始化为 nums[0]，否则全负数组会错。关键代码行必须能逐句翻译成“旧状态怎样变成新状态”，否则就只是模板。

\textbf{正确性。} 任意以 i 结尾的连续段只有两类：只含 nums[i]，或由某个以 i-1 结尾的连续段加上 nums[i]。取前一类最优即可覆盖全部候选。

\textbf{边界实验。} 先用全负数组、单元素、和溢出、答案不能初始化为零做反例检查。实验时覆盖全负数组、单元素、和溢出、答案不能初始化为零，先打印完整 DP 表，再比较空间压缩版本。若压缩版不同，优先检查遍历顺序和旧值保存。

\textbf{复杂度变化。} 暴力 O(\ensuremath{n^{2}}) 或 O(\ensuremath{n^{3}})，Kadane 算法 O(n) 时间、O(1) 空间。

\textbf{迁移追问。} 如果题目改成“若要求返回区间下标，同步记录重新开始位置和最优左右端”，先判断原状态是否仍然覆盖新答案集合，再决定是微调边界、改 value 语义，还是换算法。

\subsection{LeetCode 304 二维区域和检索 - 矩阵不可变}

\textbf{题目说明。} LeetCode 304 二维区域和检索 - 矩阵不可变给定一个固定矩阵，先构造 NumMatrix，再多次调用 sumRegion(row1, col1, row2, col2) 查询闭矩形区域和。

\textbf{样例入口。} 题面样例的三次查询分别返回 8、11、12；正文只保留查询结果，完整矩阵见原题样例。

\textbf{暴力基线。} 暴力每次查询都双层循环扫描 row1..row2、col1..col2，把矩形内元素相加。

\textbf{暴力过程。} 若多次查询同一个固定矩阵，sumRegion(2, 1, 4, 3) 会扫描一块矩形，下一次相邻矩形又会重复扫描大量格子。

\textbf{重复工作。} 题面模型：固定矩阵上需要多次查询闭矩形区域和。暴力版的慢点要落到本题：浪费在矩阵不变却每次重新读格子。任意查询矩形都可以由几个左上前缀矩形组合出来。后面的优化只消掉这类重复工作，不能改变答案集合。

\textbf{优化条件。} 本题的关键观察是：预处理二维前缀和后，每次查询用四个角做容斥。这句话必须能翻译成可维护状态；如果题目条件改变后观察不再成立，就不能继续套原来的写法。

\textbf{相邻状态。} 规律是查询公式对应“右下、上方、左侧、左上重叠”四个量；多开一行一列哨兵可以统一第一行和第一列。把这个特例继续拆开看：先标出已经确认的前缀或后缀，再标出未知区；能被丢掉的候选，必须是以后再也不会改变已确认区域语义的候选。这样才算从例子推出数组不变量，而不是只记一次操作。

\textbf{状态复用。} 预处理 prefix[r+1][c+1] 表示原矩阵 [0..r][0..c] 的和。查询时用右下大矩形减上方、减左方、加左上重叠。手算时只改被新输入影响的那一小部分，而不是回头重算完整候选。

\textbf{指针和字段。} 状态字段要能说清：prefix 的多一行多一列是哨兵，row1/col1/row2/col2 是原矩阵闭区间坐标，查询公式使用加一后的前缀坐标。手算时按这个协议跟踪：sumRegion(2, 1, 4, 3) 等于 P[5][4]-P[2][4]-P[5][1]+P[2][1]。

\textbf{代码落点。} 构建时分四项：当前格、上方前缀、左方前缀，再减左上重叠。代码里分别对应 matrix 当前值、上方 prefix、左方 prefix 和左上 prefix。关键代码行必须能逐句翻译成“旧状态怎样变成新状态”，否则就只是模板。

\textbf{正确性。} 右下大矩形包含目标、上方和左方；上方左方相减时左上角被减了两次，所以要加回一次。

\textbf{边界实验。} 先用第一行第一列、单行单列、空矩阵、多次重复查询、坐标加一做反例检查。实验时把第一行第一列、单行单列、空矩阵、多次重复查询、坐标加一列成表，再打印关键下标和有效区间。若某个元素被写入后又被未来步骤破坏，说明区间不变量没有成立。

\textbf{复杂度变化。} 单次暴力 O(矩形面积)，预处理 O(mn)，每次查询 O(1)，空间 O(mn)。

\textbf{迁移追问。} 如果题目改成“若矩阵会更新，普通二维前缀和失效，需要二维树状数组或线段树”，先判断原状态是否仍然覆盖新答案集合，再决定是微调边界、改 value 语义，还是换算法。

\subsection{LeetCode 238 除自身以外数组的乘积}

\textbf{题目说明。} LeetCode 238 除自身以外数组的乘积给定整数数组 nums，返回 answer，其中 answer[i] 等于 nums 中除 nums[i] 以外所有元素的乘积，且不能使用除法。

\textbf{样例入口。} 样例 nums=[1, 2, 3, 4]，输出 [24, 12, 8, 6]；每个位置都排除自身后乘其余元素。

\textbf{暴力基线。} 暴力对每个位置 i 扫描整个数组，跳过 i，把其他元素全部相乘得到 answer[i]。

\textbf{暴力过程。} 以 [1, 2, 3, 4] 为例，answer[2] 要算 1*2*4；answer[3] 又重算 1*2*3，左侧乘积被重复计算。

\textbf{重复工作。} 题面模型：给定整数数组，返回 answer，其中 answer[i] 等于 nums[i] 之外所有元素的乘积。暴力版的慢点要落到本题：浪费在每个 i 都重算左侧和右侧乘积。位置 i 的答案天然等于 i 左侧乘积乘 i 右侧乘积。后面的优化只消掉这类重复工作，不能改变答案集合。

\textbf{优化条件。} 本题的关键观察是：不能用除法时，答案由左侧乘积和右侧乘积组成；先写入左侧前缀积，再从右向左乘右侧后缀积。这句话必须能翻译成可维护状态；如果题目条件改变后观察不再成立，就不能继续套原来的写法。

\textbf{相邻状态。} 规律是答案等于左侧乘积乘右侧乘积，不能使用除法时就用两次扫描合并贡献。把这个特例继续拆开看：先标出已经确认的前缀或后缀，再标出未知区；能被丢掉的候选，必须是以后再也不会改变已确认区域语义的候选。这样才算从例子推出数组不变量，而不是只记一次操作。

\textbf{状态复用。} 第一遍从左到右，把 answer[i] 写成 i 左侧所有数的乘积；第二遍从右到左，用 right\_product 乘回右侧贡献。手算时只改被新输入影响的那一小部分，而不是回头重算完整候选。

\textbf{指针和字段。} 状态字段要能说清：left\_product 隐含在 answer[i]，right\_product 是当前 i 右侧乘积，answer[i] 最终是左右贡献相乘。手算时按这个协议跟踪：[1, 2, 3, 4] 第一遍 answer=[1, 1, 2, 6]；第二遍 right\_product 依次为 1、4、12、24，最终 [24, 12, 8, 6]。

\textbf{代码落点。} 关键是第二遍先 answer[i]*=right\_product，再 right\_product*=nums[i]，这样不会把 nums[i] 自己乘进去。关键代码行必须能逐句翻译成“旧状态怎样变成新状态”，否则就只是模板。

\textbf{正确性。} 第一遍保证 answer[i] 只含左侧元素，第二遍乘入的 right\_product 只含右侧元素，二者乘积正好排除自身。

\textbf{边界实验。} 先用一个零、多个零、负数、乘积溢出做反例检查。实验时把一个零、多个零、负数、乘积溢出列成表，再打印关键下标和有效区间。若某个元素被写入后又被未来步骤破坏，说明区间不变量没有成立。

\textbf{复杂度变化。} 暴力 O(\ensuremath{n^{2}})，前后缀合并 O(n) 时间；若输出数组不计额外空间，则额外空间 O(1)。

\textbf{迁移追问。} 如果题目改成“这个题展示输出数组可作为状态存储的一部分”，先判断原状态是否仍然覆盖新答案集合，再决定是微调边界、改 value 语义，还是换算法。

\subsection{LeetCode 11 盛最多水的容器}

\textbf{题目说明。} LeetCode 11 盛最多水的容器给定数组 height，height[i] 表示第 i 条竖线高度。选择两条线与 x 轴形成容器，返回能容纳的最大水量。

\textbf{样例入口。} 样例 height=[1, 8, 6, 2, 5, 4, 8, 3, 7]，输出 49；选择下标 1 和 8，宽度 7，短板高度 7，面积 49。

\textbf{暴力基线。} 暴力枚举所有 left<right 的下标对，面积等于 (right-left)*min(height[left], height[right])，取最大值。

\textbf{暴力过程。} 以 [1, 8, 6, 2, 5, 4, 8, 3, 7] 为例，暴力会试 (0, 1)、(0, 2) 一直到 (0, 8)，再试 (1, 2)、(1, 3)。大量候选只是宽度和短板变化。

\textbf{重复工作。} 题面模型：给定每个下标的高度，选择两条竖线和横轴组成容器，返回能盛水的最大面积。暴力版的慢点要落到本题：浪费在逐个试所有右端。若当前短板是 left，保留 left 再把 right 左移只会让宽度变小，短板高度不可能超过 height[left]，这些候选都被当前状态支配。后面的优化只消掉这类重复工作，不能改变答案集合。

\textbf{优化条件。} 本题的关键观察是：每次移动短板，因为移动长板只会缩小宽度且不会提高短板。这句话必须能翻译成可维护状态；如果题目条件改变后观察不再成立，就不能继续套原来的写法。

\textbf{相邻状态。} 规律是每轮淘汰短板那一侧，因为所有保留它且宽度更小的候选都被当前面积上界支配。把这个特例继续拆开看：右端进入只带来一个新元素，左端离开只删掉一个旧元素；只有当旧左端被移走后不可能再产生更优答案时，窗口才可以单向推进。若这个单调淘汰条件不存在，就要换成前缀和、队列或其他状态。

\textbf{状态复用。} 双指针从两端开始，维护 left、right 和 best。每轮计算当前面积，然后移动高度较小的一侧。手算时只改被新输入影响的那一小部分，而不是回头重算完整候选。

\textbf{指针和字段。} 状态字段要能说清：left/right 是当前最宽候选的两端，shorter 是限制高度的一侧，best 是已见过最大面积。手算时按这个协议跟踪：样例中 left=0 高 1、right=8 高 7，面积 8。短板是 left，移动 left 到 1；此时高 8 和 7，面积 49，更新答案。

\textbf{代码落点。} 关键代码是 if (height[left] < height[right]) ++left; else --right; 。移动长板无法提高短板上界，所以不移动长板。关键代码行必须能逐句翻译成“旧状态怎样变成新状态”，否则就只是模板。

\textbf{正确性。} 每次移动短板时，所有包含这个短板且宽度更小的候选面积都不可能超过当前面积上界，因此可以一次排除一批候选。

\textbf{边界实验。} 先用两端高度相等、数组长度为二、面积乘法可能溢出做反例检查。实验时重点跑两端高度相等、数组长度为二、面积乘法可能溢出，并打印左右端、窗口计数和答案。若左端发生回退，或者窗口失去合法性后仍更新答案，就能很快定位错误。

\textbf{复杂度变化。} 暴力 O(\ensuremath{n^{2}})，双指针每轮移动一端，总共 O(n) 轮，时间 O(n)，空间 O(1)。

\textbf{迁移追问。} 如果题目改成“接雨水计算每个位置贡献，不能把本题双指针证明直接搬过去”，先判断原状态是否仍然覆盖新答案集合，再决定是微调边界、改 value 语义，还是换算法。

\subsection{LeetCode 167 两数之和 II 输入有序数组}

\textbf{题目说明。} LeetCode 167 两数之和 II 输入有序数组 的题面入口要先还原成具体任务：有序数组中左右指针的和随指针移动单调变化。

\textbf{样例入口。} 拿 [2, 7, 11, 15]、target=9 手算，左右和为 17 太大，右指针左移；2+7 正好等于 9。

\textbf{暴力基线。} 暴力枚举所有 i<j 的下标对，检查 numbers[i]+numbers[j] 是否等于 target。

\textbf{暴力过程。} 以 [2, 7, 11, 15]、target=9 为例，暴力会试 2+7、2+11、2+15...；但数组有序，2+15 太大时，右端继续更大没有意义。

\textbf{重复工作。} 题面模型：有序数组中左右指针的和随指针移动单调变化。暴力版的慢点要落到本题：浪费在没有利用有序性。当前和太小，固定右端时更小的左端都不可能达标；当前和太大，固定左端时更大的右端都不可能达标。后面的优化只消掉这类重复工作，不能改变答案集合。

\textbf{优化条件。} 本题的关键观察是：当前和太小左指针右移，太大右指针左移，等于目标返回。这句话必须能翻译成可维护状态；如果题目条件改变后观察不再成立，就不能继续套原来的写法。

\textbf{相邻状态。} 规律是有序数组让和随左指针右移变大、随右指针左移变小，因此每次能排除一侧候选。把这个特例继续拆开看：右端进入只带来一个新元素，左端离开只删掉一个旧元素；只有当旧左端被移走后不可能再产生更优答案时，窗口才可以单向推进。若这个单调淘汰条件不存在，就要换成前缀和、队列或其他状态。

\textbf{状态复用。} left 从 0 开始，right 从 n-1 开始。每轮计算 sum，sum 小则 left++，sum 大则 right--，相等返回 1 基下标。手算时只改被新输入影响的那一小部分，而不是回头重算完整候选。

\textbf{指针和字段。} 状态字段要能说清：left/right 是当前候选对，sum 是当前两数和，target 决定排除哪一侧。手算时按这个协议跟踪：[2, 7, 11, 15] 中 2+15=17 太大，right 左移；2+11=13 仍大，再左移；2+7=9 返回 [1, 2]。

\textbf{代码落点。} 关键代码是 sum<target 时 ++left，sum>target 时 --right。返回时要加 1，因为题目要求 1 基下标。关键代码行必须能逐句翻译成“旧状态怎样变成新状态”，否则就只是模板。

\textbf{正确性。} 有序性保证移动一侧能排除一整批不可能候选，且剩余区间仍包含可能答案。

\textbf{边界实验。} 先用两个数在边界、重复值、题目下标从一开始、保证有解做反例检查。实验时重点跑两个数在边界、重复值、题目下标从一开始、保证有解，并打印左右端、窗口计数和答案。若左端发生回退，或者窗口失去合法性后仍更新答案，就能很快定位错误。

\textbf{复杂度变化。} 暴力 O(\ensuremath{n^{2}})，双指针 O(n) 时间、O(1) 空间。

\textbf{迁移追问。} 如果题目改成“无序版本通常用哈希表，不能直接套双指针”，先判断原状态是否仍然覆盖新答案集合，再决定是微调边界、改 value 语义，还是换算法。

\subsection{LeetCode 76 最小覆盖子串}

\textbf{题目说明。} LeetCode 76 最小覆盖子串给定字符串 s 和 t，要求在 s 中找最短连续子串，使它包含 t 中每个字符及其次数；不存在则返回空串。

\textbf{样例入口。} 样例 s=ADOBECODEBANC，t=ABC，输出 BANC；它是覆盖 A、B、C 的最短连续子串。

\textbf{暴力基线。} 暴力枚举 s 的每个子串，重新统计字符频次，检查是否覆盖 t 的每个字符及其次数。

\textbf{暴力过程。} 以 s=ADOBECODEBANC、t=ABC 为例，暴力会检查 A、AD、ADO...直到 ADOBEC 才第一次覆盖；下一轮从 D 开始又重新统计 DOBEC。

\textbf{重复工作。} 题面模型：给定字符串 s 和 t，在 s 中找一个最短连续子串，使它覆盖 t 的每个字符及其次数。暴力版的慢点要落到本题：浪费在相邻子串重复统计字符。右端扩展只新增一个字符，左端收缩只移走一个字符。后面的优化只消掉这类重复工作，不能改变答案集合。

\textbf{优化条件。} 本题的关键观察是：窗口内字符计数必须覆盖目标计数，满足后尽量收缩左端；用 formed 记录满足需求的字符种类，避免每次全量比较。这句话必须能翻译成可维护状态；如果题目条件改变后观察不再成立，就不能继续套原来的写法。

\textbf{相邻状态。} 规律是右端负责获得可行，左端负责在可行时尽量缩短；计数表要区分已有字符和需求字符。把这个特例继续拆开看：右端进入只带来一个新元素，左端离开只删掉一个旧元素；只有当旧左端被移走后不可能再产生更优答案时，窗口才可以单向推进。若这个单调淘汰条件不存在，就要换成前缀和、队列或其他状态。

\textbf{状态复用。} need 保存目标频次，window 保存当前窗口频次，formed 或 required\_count 记录当前满足需求的字符种类；满足后尽量收缩 left。手算时只改被新输入影响的那一小部分，而不是回头重算完整候选。

\textbf{指针和字段。} 状态字段要能说清：right 负责扩展获得覆盖，left 负责在覆盖时缩短窗口，best\_left/best\_len 记录最短答案。手算时按这个协议跟踪：ADOBEC 第一次覆盖 ABC，记录长度 6；左端移走 A 后不覆盖，停止收缩。继续扩展到 BANC 时再次覆盖，收缩得到长度 4。

\textbf{代码落点。} 关键顺序是加入 s[right] 后更新 formed；while formed==required 时先更新答案，再移走 s[left] 并可能降低 formed。关键代码行必须能逐句翻译成“旧状态怎样变成新状态”，否则就只是模板。

\textbf{正确性。} 每个 right 下，while 循环会把 left 推到刚好不能再缩的位置，因此不会错过以该 right 结尾的最短覆盖窗口。

\textbf{边界实验。} 先用目标有重复字符、源串缺字符、最小窗口在开头或结尾、大小写做反例检查。实验时重点跑目标有重复字符、源串缺字符、最小窗口在开头或结尾、大小写，并打印左右端、窗口计数和答案。若左端发生回退，或者窗口失去合法性后仍更新答案，就能很快定位错误。

\textbf{复杂度变化。} 暴力至少 O(\ensuremath{n^{2}}*字符集)，滑动窗口每个字符进出一次，时间 O(n+字符集)，空间 O(字符集)。

\textbf{迁移追问。} 如果题目改成“若字符集固定，数组计数更快；若是 Unicode，哈希表更通用”，先判断原状态是否仍然覆盖新答案集合，再决定是微调边界、改 value 语义，还是换算法。

\subsection{LeetCode 209 长度最小的子数组}

\textbf{题目说明。} LeetCode 209 长度最小的子数组给定正整数数组 nums 和目标 target，返回和至少为 target 的最短连续子数组长度；不存在则返回 0。

\textbf{样例入口。} 样例 target=7，nums=[2, 3, 1, 2, 4, 3]，输出 2；最短合法连续段是 [4, 3]。

\textbf{暴力基线。} 暴力枚举每个起点 left，向右累计 sum，第一次 sum>=target 时更新长度，再换下一个起点。

\textbf{暴力过程。} 以 target=7、nums=[2, 3, 1, 2, 4, 3] 为例，从 0 开始要扩到 [2, 3, 1, 2]；从 1 开始又重新累加 [3, 1, 2, 4]。

\textbf{重复工作。} 题面模型：给定正整数数组和目标值，返回和至少为目标值的最短连续子数组长度。暴力版的慢点要落到本题：浪费在相邻起点重复累加。因为 nums 全为正数，右端增加只会让 sum 变大，左端移走只会让 sum 变小。后面的优化只消掉这类重复工作，不能改变答案集合。

\textbf{优化条件。} 本题的关键观察是：正数数组让窗口和随右端增加、随左端减少；右端扩展到满足目标后，不断左移缩短并更新答案。这句话必须能翻译成可维护状态；如果题目条件改变后观察不再成立，就不能继续套原来的写法。

\textbf{相邻状态。} 规律是全正数保证右扩和变大、左缩和变小，所以每个左端只需被移动一次。把这个特例继续拆开看：右端进入只带来一个新元素，左端离开只删掉一个旧元素；只有当旧左端被移走后不可能再产生更优答案时，窗口才可以单向推进。若这个单调淘汰条件不存在，就要换成前缀和、队列或其他状态。

\textbf{状态复用。} left、right 维护当前窗口和 sum。right 每次加入一个数；当 sum>=target 时，循环收缩 left 并更新最短长度。手算时只改被新输入影响的那一小部分，而不是回头重算完整候选。

\textbf{指针和字段。} 状态字段要能说清：sum 是 nums[left..right] 的和，answer 是已见过最短合法长度；left 只在窗口合法时右移。手算时按这个协议跟踪：样例中 right 到下标 3 时 sum=8，记录长度 4，移走 2 后 sum=6 停止；后来窗口 [4, 3] sum=7，答案更新为 2。

\textbf{代码落点。} 关键代码是 sum += nums[right]; while (sum >= target) \{ ans=min(ans, right-left+1); sum-=nums[left++]; \}。关键代码行必须能逐句翻译成“旧状态怎样变成新状态”，否则就只是模板。

\textbf{正确性。} 正数保证收缩 left 后 sum 单调下降；对每个 right，while 会找到以它结尾的所有可缩短合法窗口，不会漏最短值。

\textbf{边界实验。} 先用不存在答案、单元素达到目标、全正是必要条件做反例检查。实验时重点跑不存在答案、单元素达到目标、全正是必要条件，并打印左右端、窗口计数和答案。若左端发生回退，或者窗口失去合法性后仍更新答案，就能很快定位错误。

\textbf{复杂度变化。} 暴力 O(\ensuremath{n^{2}})，窗口中每个元素至多进入一次、离开一次，时间 O(n)，空间 O(1)。

\textbf{迁移追问。} 如果题目改成“若含负数，需要前缀和加单调队列”，先判断原状态是否仍然覆盖新答案集合，再决定是微调边界、改 value 语义，还是换算法。

\subsection{LeetCode 283 移动零}

\textbf{题目说明。} LeetCode 283 移动零给定数组 nums，要求原地把所有 0 移到末尾，同时保持非零元素的相对顺序。

\textbf{样例入口。} 样例 nums=[0, 1, 0, 3, 12]，处理后为 [1, 3, 12, 0, 0]；非零元素 1、3、12 的相对顺序不变。

\textbf{暴力基线。} 暴力可以新建数组，先放所有非零元素，再补同样数量的零，最后复制回原数组。

\textbf{暴力过程。} 以 [0, 1, 0, 3, 12] 为例，辅助数组先收集 [1, 3, 12]，再补 [0, 0]。

\textbf{重复工作。} 题面模型：给定数组，原地把所有 0 移到末尾，同时保持非零元素的相对顺序。暴力版的慢点要落到本题：浪费在辅助数组。非零元素的相对顺序由读取顺序天然保证，只需要一个原地写入位置。后面的优化只消掉这类重复工作，不能改变答案集合。

\textbf{优化条件。} 本题的关键观察是：稳定保留非零元素顺序，把零集中到尾部；慢指针写入下一个非零位置，最后补零或用交换。这句话必须能翻译成可维护状态；如果题目条件改变后观察不再成立，就不能继续套原来的写法。

\textbf{相邻状态。} 规律是先稳定压缩非零前缀，再填零；不能在扫描时盲目交换导致顺序变化。把这个特例继续拆开看：先标出已经确认的前缀或后缀，再标出未知区；能被丢掉的候选，必须是以后再也不会改变已确认区域语义的候选。这样才算从例子推出数组不变量，而不是只记一次操作。

\textbf{状态复用。} write 指向下一个非零元素应写入的位置。第一遍把非零元素稳定写到前面，第二遍把 write 之后填 0。手算时只改被新输入影响的那一小部分，而不是回头重算完整候选。

\textbf{指针和字段。} 状态字段要能说清：read 是扫描位置，write 左侧是不含 0 的稳定前缀，write 到末尾最后会被置零。手算时按这个协议跟踪：读到 1 写到 nums[0]，读到 3 写到 nums[1]，读到 12 写到 nums[2]，最后 nums[3]、nums[4] 置 0。

\textbf{代码落点。} 关键代码是 if (nums[read]!=0) nums[write++]=nums[read]; 后面 while(write<n) nums[write++]=0。关键代码行必须能逐句翻译成“旧状态怎样变成新状态”，否则就只是模板。

\textbf{正确性。} 所有非零元素按原读取顺序写入前缀，因此稳定性保持；剩余位置数量等于原零的数量，补零后满足题意。

\textbf{边界实验。} 先用全零、无零、零在开头结尾、稳定顺序要求做反例检查。实验时把全零、无零、零在开头结尾、稳定顺序要求列成表，再打印关键下标和有效区间。若某个元素被写入后又被未来步骤破坏，说明区间不变量没有成立。

\textbf{复杂度变化。} 辅助数组 O(n) 空间；原地稳定写 O(n) 时间、O(1) 额外空间。

\textbf{迁移追问。} 如果题目改成“若目标是移动指定值，模型与移除元素相同”，先判断原状态是否仍然覆盖新答案集合，再决定是微调边界、改 value 语义，还是换算法。

\subsection{LeetCode 189 轮转数组}

\textbf{题目说明。} LeetCode 189 轮转数组给定数组 nums 和整数 k，要求把数组原地向右轮转 k 个位置。

\textbf{样例入口。} 样例 nums=[1, 2, 3, 4, 5, 6, 7]，k=3，轮转后为 [5, 6, 7, 1, 2, 3, 4]。

\textbf{暴力基线。} 暴力可以新建 answer，让 answer[(i+k)\%n]=nums[i]，最后复制回 nums。

\textbf{暴力过程。} 以 [1, 2, 3, 4, 5, 6, 7]、k=3 为例，1 去下标 3，5 去下标 0，结果是 [5, 6, 7, 1, 2, 3, 4]。

\textbf{重复工作。} 题面模型：给定数组和整数 k，将数组中的元素整体向右轮转 k 个位置。暴力版的慢点要落到本题：浪费在额外数组。右移 k 位等价于把后 k 个元素整体搬到前面，原地可以用反转实现两段换序。后面的优化只消掉这类重复工作，不能改变答案集合。

\textbf{优化条件。} 本题的关键观察是：右移 k 位等价于把数组切成两段后换序，三次反转能原地完成：整体、前 k、后 n 减 k。这句话必须能翻译成可维护状态；如果题目条件改变后观察不再成立，就不能继续套原来的写法。

\textbf{相邻状态。} 规律是右移 k 等价于两段换序，k 要先对 n 取模。把这个特例继续拆开看：先标出已经确认的前缀或后缀，再标出未知区；能被丢掉的候选，必须是以后再也不会改变已确认区域语义的候选。这样才算从例子推出数组不变量，而不是只记一次操作。

\textbf{状态复用。} 先 k\%=n；整体反转让后段到前面，再反转前 k 段和后 n-k 段，恢复每段内部顺序。手算时只改被新输入影响的那一小部分，而不是回头重算完整候选。

\textbf{指针和字段。} 状态字段要能说清：k 是有效轮转距离，三个区间分别是 whole、[0, k)、[k, n)。手算时按这个协议跟踪：[1, 2, 3, 4, 5, 6, 7] 整体反转为 [7, 6, 5, 4, 3, 2, 1]；前 3 个反转为 [5, 6, 7, 4, 3, 2, 1]；后 4 个反转为 [5, 6, 7, 1, 2, 3, 4]。

\textbf{代码落点。} 关键是空数组先返回，否则 k 取模会除零。随后按“整体反转、修正前段、修正后段”三步写，不把长函数调用塞成一整行。关键代码行必须能逐句翻译成“旧状态怎样变成新状态”，否则就只是模板。

\textbf{正确性。} 整体反转完成两段位置交换但段内顺序反了；再分别反转两段，正好得到右移后的段顺序。

\textbf{边界实验。} 先用k 为零、k 大于 n、空数组、单元素做反例检查。实验时把k 为零、k 大于 n、空数组、单元素列成表，再打印关键下标和有效区间。若某个元素被写入后又被未来步骤破坏，说明区间不变量没有成立。

\textbf{复杂度变化。} 辅助数组 O(n) 空间；三次反转 O(n) 时间、O(1) 额外空间。

\textbf{迁移追问。} 如果题目改成“环状替换也可做，但要处理最大公约数个循环”，先判断原状态是否仍然覆盖新答案集合，再决定是微调边界、改 value 语义，还是换算法。

\subsection{LeetCode 27 移除元素}

\textbf{题目说明。} LeetCode 27 移除元素 的题面入口要先还原成具体任务：原地过滤，慢指针左侧始终是不等于目标值的稳定前缀。

\textbf{样例入口。} 拿 [3, 2, 2, 3]、val=3 手算，稳定写法读到 3 不写，读到两个 2 依次写入前缀，最后有效区间是 [2, 2]。若题目不要求顺序，也可以把左侧 3 和右侧非 3 交换来减少写入。

\textbf{暴力基线。} 暴力可以新建数组，只复制不等于 val 的元素；最后返回新数组长度或把结果写回原数组前缀。

\textbf{暴力过程。} 以 nums=[3, 2, 2, 3]、val=3 为例，辅助数组依次跳过 3、写入 2、写入 2、跳过 3，得到有效前缀 [2, 2]。

\textbf{重复工作。} 题面模型：原地过滤，慢指针左侧始终是不等于目标值的稳定前缀。暴力版的慢点要落到本题：浪费在辅助数组。若只需要原地返回有效长度，辅助数组的写入位置就是原数组里的 write 指针。后面的优化只消掉这类重复工作，不能改变答案集合。

\textbf{优化条件。} 本题的关键观察是：保持顺序用快慢指针；不保持顺序可用左右交换减少写入。这句话必须能翻译成可维护状态；如果题目条件改变后观察不再成立，就不能继续套原来的写法。

\textbf{相邻状态。} 规律是先确认是否稳定，再决定快慢指针还是左右交换；所有元素都被移除时返回长度 0。把这个特例继续拆开看：先标出已经确认的前缀或后缀，再标出未知区；能被丢掉的候选，必须是以后再也不会改变已确认区域语义的候选。这样才算从例子推出数组不变量，而不是只记一次操作。

\textbf{状态复用。} 稳定写法维护 write，nums[0..write-1] 都是不等于 val 的元素且相对顺序不变；read 扫描所有元素。手算时只改被新输入影响的那一小部分，而不是回头重算完整候选。

\textbf{指针和字段。} 状态字段要能说清：read 是输入扫描位置，write 是下一个保留元素写入位置，返回值 write 是有效长度，不是数组物理容量。手算时按这个协议跟踪：读到第一个 3 不写；读到 2 写到 nums[0]；读到第二个 2 写到 nums[1]；最后 3 不写，返回 2。

\textbf{代码落点。} 关键代码是 if (nums[read] != val) nums[write++] = nums[read]; 。若题目明确不要求稳定顺序，才考虑左右交换版本。关键代码行必须能逐句翻译成“旧状态怎样变成新状态”，否则就只是模板。

\textbf{正确性。} write 左侧只接收读过且不等于 val 的元素，read 每前进一步都保持这个前缀语义，扫描结束时前缀正好包含所有应保留元素。

\textbf{边界实验。} 先用所有元素都被移除、没有元素被移除、目标值在尾部密集做反例检查。实验时把所有元素都被移除、没有元素被移除、目标值在尾部密集列成表，再打印关键下标和有效区间。若某个元素被写入后又被未来步骤破坏，说明区间不变量没有成立。

\textbf{复杂度变化。} 辅助数组 O(n) 时间和 O(n) 空间；原地快慢指针 O(n) 时间、O(1) 额外空间。

\textbf{迁移追问。} 如果题目改成“工程里要先确认是否要求稳定顺序，不同要求对应不同写法”，先判断原状态是否仍然覆盖新答案集合，再决定是微调边界、改 value 语义，还是换算法。

\subsection{LeetCode 15 三数之和}

\textbf{题目说明。} LeetCode 15 三数之和给定整数数组 nums，要求找出所有不重复三元组 [a, b, c]，使 a+b+c=0。

\textbf{样例入口。} 样例 nums=[-1, 0, 1, 2, -1, -4]，输出 [[-1, -1, 2], [-1, 0, 1]]；结果里不能出现重复三元组。

\textbf{暴力基线。} 暴力三层循环枚举 i、j、k，检查 nums[i]+nums[j]+nums[k] 是否为 0，再用集合去掉重复三元组。

\textbf{暴力过程。} 以 [-1, 0, 1, 2, -1, -4] 为例，暴力会枚举 (-1, 0, 1)、(-1, 0, 2)、(-1, 0, -1) 等所有三元组；两个 -1 会制造重复答案。

\textbf{重复工作。} 题面模型：给定整数数组，返回所有不重复三元组，使三个数之和等于零。暴力版的慢点要落到本题：浪费在第三层枚举和事后去重。排序后固定一个数，剩余两个数变成有序两数之和，可以根据当前和一次排除一侧。后面的优化只消掉这类重复工作，不能改变答案集合。

\textbf{优化条件。} 本题的关键观察是：排序后固定第一个数，剩余两数转成有序双指针；固定值和左右值都要跳过重复。这句话必须能翻译成可维护状态；如果题目条件改变后观察不再成立，就不能继续套原来的写法。

\textbf{相邻状态。} 规律是排序把三数问题拆成枚举固定点加双指针，并在每层处理去重。把这个特例继续拆开看：右端进入只带来一个新元素，左端离开只删掉一个旧元素；只有当旧左端被移走后不可能再产生更优答案时，窗口才可以单向推进。若这个单调淘汰条件不存在，就要换成前缀和、队列或其他状态。

\textbf{状态复用。} 先排序。外层 i 固定第一个数，内层 left=i+1、right=n-1 搜索 target=-nums[i]；找到答案后跳过相同 left/right 值。手算时只改被新输入影响的那一小部分，而不是回头重算完整候选。

\textbf{指针和字段。} 状态字段要能说清：i 是本层固定点，left/right 是右侧有序区间的候选两端，sum 是三数和，res 保存已去重结果。手算时按这个协议跟踪：排序后 [-4, -1, -1, 0, 1, 2]。i=1 固定 -1，left 指向 -1，right 指向 2，和为 0，记录 [-1, -1, 2]，然后跳过重复。

\textbf{代码落点。} 关键代码是 sum<0 时 ++left，sum>0 时 --right；i>0 且 nums[i]==nums[i-1] 要 continue，命中后 left/right 也要跳过重复值。关键代码行必须能逐句翻译成“旧状态怎样变成新状态”，否则就只是模板。

\textbf{正确性。} 排序后，当前和太小，固定 i 和 right 时更小的 left 都不可能达标；当前和太大，固定 i 和 left 时更大的 right 都不可能达标。去重发生在同一层，避免漏掉不同层组合。

\textbf{边界实验。} 先用全零、重复元素很多、没有答案、结果顺序不要求但组合不能重复做反例检查。实验时重点跑全零、重复元素很多、没有答案、结果顺序不要求但组合不能重复，并打印左右端、窗口计数和答案。若左端发生回退，或者窗口失去合法性后仍更新答案，就能很快定位错误。

\textbf{复杂度变化。} 暴力 O(\ensuremath{n^{3}}) 加去重开销；排序 O(n log n)，外层 n 次、内层双指针线性，总时间 O(\ensuremath{n^{2}})，输出空间另计。

\textbf{迁移追问。} 如果题目改成“四数之和再外层固定一个数，并用 long long 防止目标和溢出”，先判断原状态是否仍然覆盖新答案集合，再决定是微调边界、改 value 语义，还是换算法。

\subsection{LeetCode 18 四数之和}

\textbf{题目说明。} LeetCode 18 四数之和 的题面入口要先还原成具体任务：两层固定加一层双指针，关键是剪枝和去重。

\textbf{样例入口。} 拿 [1, 0, -1, 0, -2, 2]、target=0 手算，排序后先固定 -2，再固定 -1，剩余两数用双指针找 3。找到 [-2, -1, 1, 2] 后，左右指针相同值必须跳过，否则会重复输出。

\textbf{暴力基线。} 暴力四层循环枚举四个下标，检查和是否等于 target，再用集合去重。

\textbf{暴力过程。} 以 [1, 0, -1, 0, -2, 2]、target=0 为例，暴力会试所有四元组，[-2, -1, 1, 2] 和含重复 0 的组合都要事后去重。

\textbf{重复工作。} 题面模型：两层固定加一层双指针，关键是剪枝和去重。暴力版的慢点要落到本题：浪费在后两层仍然盲枚举。排序后两层固定，剩余两数是有序配对问题；同层重复值可以直接跳过。后面的优化只消掉这类重复工作，不能改变答案集合。

\textbf{优化条件。} 本题的关键观察是：排序提供单调性，外层可用最小可能和和最大可能和提前跳过。这句话必须能翻译成可维护状态；如果题目条件改变后观察不再成立，就不能继续套原来的写法。

\textbf{相邻状态。} 规律是两层固定把四数降成有序两数，外层用最小可能和与最大可能和剪枝，内外层都要去重。把这个特例继续拆开看：右端进入只带来一个新元素，左端离开只删掉一个旧元素；只有当旧左端被移走后不可能再产生更优答案时，窗口才可以单向推进。若这个单调淘汰条件不存在，就要换成前缀和、队列或其他状态。

\textbf{状态复用。} 排序后固定 i 和 j，left=j+1、right=n-1 搜索 target-nums[i]-nums[j]。外层可用最小可能和、最大可能和剪枝。手算时只改被新输入影响的那一小部分，而不是回头重算完整候选。

\textbf{指针和字段。} 状态字段要能说清：i/j 是前两个固定值，left/right 是剩余区间候选，sum 用 long long 保存，res 保存去重四元组。手算时按这个协议跟踪：排序后 [-2, -1, 0, 0, 1, 2]。i=-2、j=-1 时，left=0、right=2，四数和为 -1，left 右移；随后找到 [-2, -1, 1, 2]。

\textbf{代码落点。} 关键代码和三数之和一样，但外层有两层去重。sum 必须用 long long，避免 target 和 nums 值接近 int 边界时溢出。关键代码行必须能逐句翻译成“旧状态怎样变成新状态”，否则就只是模板。

\textbf{正确性。} 两层固定后，有序双指针的排除理由和两数之和相同；同层去重只跳过值相同且角色相同的候选，不会跳过需要使用的重复元素。

\textbf{边界实验。} 先用重复元素、目标超出 int、答案很多导致输出空间主导复杂度做反例检查。实验时重点跑重复元素、目标超出 int、答案很多导致输出空间主导复杂度，并打印左右端、窗口计数和答案。若左端发生回退，或者窗口失去合法性后仍更新答案，就能很快定位错误。

\textbf{复杂度变化。} 暴力 O(\ensuremath{n^{4}})，排序后 O(\ensuremath{n^{3}})，输出空间另计。

\textbf{迁移追问。} 如果题目改成“k 数之和可以递归降维，但工程实现要控制重复和边界”，先判断原状态是否仍然覆盖新答案集合，再决定是微调边界、改 value 语义，还是换算法。

\subsection{LeetCode 3 无重复字符的最长子串}

\textbf{题目说明。} LeetCode 3 无重复字符的最长子串 的题面入口要先还原成具体任务：窗口保存当前无重复子串，右端每次加入一个字符。

\textbf{样例入口。} 拿字符串 abca 手算。窗口 abc 合法，右端加入第二个 a 后，真正冲突的是上一次 a 的位置；左端从 0 跳到 1 后，bca 重新合法。再试 abba：处理第二个 b 时左端跳到 2，后面遇到 a 的上次位置在窗口左侧之外，不能让 left 回退。

\textbf{暴力基线。} 暴力固定每个起点 left，向右逐个加入字符，用集合检查是否重复；一旦重复就停止当前起点，换下一个起点重新建集合。

\textbf{暴力过程。} 以 s=abcabcbb 为例，暴力会检查 a、ab、abc，到 abca 时发现 a 重复；再从 b 重新检查 b、bc、bca。相邻起点的窗口 b、c 其实已经在上一轮出现过。

\textbf{重复工作。} 题面模型：窗口保存当前无重复子串，右端每次加入一个字符。暴力版的慢点要落到本题：浪费在每换一个 left 就重新统计一段已经看过的字符。真正变化只是左端离开若干旧字符，右端继续加入一个新字符。后面的优化只消掉这类重复工作，不能改变答案集合。

\textbf{优化条件。} 本题的关键观察是：用上次出现位置直接跳过无效左端，左边界只能向右不能回退。这句话必须能翻译成可维护状态；如果题目条件改变后观察不再成立，就不能继续套原来的写法。

\textbf{相邻状态。} 规律是左边界只向右，状态保存每个字符最近出现位置。把这个特例继续拆开看：右端进入只带来一个新元素，左端离开只删掉一个旧元素；只有当旧左端被移走后不可能再产生更优答案时，窗口才可以单向推进。若这个单调淘汰条件不存在，就要换成前缀和、队列或其他状态。

\textbf{状态复用。} 保存窗口左端 left、右端 right、每个字符最近出现位置 last、当前最长答案 best。遇到重复字符时，left 直接跳到 last[ch]+1，但只能向右跳。手算时只改被新输入影响的那一小部分，而不是回头重算完整候选。

\textbf{指针和字段。} 状态字段要能说清：right 是当前加入字符位置，left 是当前无重复窗口起点，last[ch] 是字符 ch 上一次出现下标，best 是已经见过的最大合法窗口长度。手算时按这个协议跟踪：手算 abba：right=0 加 a，left=0；right=1 加 b；right=2 再遇 b，left 跳到 2；right=3 遇 a，但上次 a 在 0，已经在窗口外，left 不能回退。

\textbf{代码落点。} 关键代码是 left=max(left, last[ch]+1)，然后 last[ch]=right，最后用 right-left+1 更新答案。顺序不能反，否则会把当前字符当成自己的历史位置。关键代码行必须能逐句翻译成“旧状态怎样变成新状态”，否则就只是模板。

\textbf{正确性。} left 跳过的是所有仍包含重复 ch 的起点，这些起点不可能形成以 right 结尾的合法窗口；left 不回退保证每个字符至多进出窗口一次。

\textbf{边界实验。} 先用空串、全相同字符、重复字符出现在窗口左侧之外、扩展字符集做反例检查。实验时重点跑空串、全相同字符、重复字符出现在窗口左侧之外、扩展字符集，并打印左右端、窗口计数和答案。若左端发生回退，或者窗口失去合法性后仍更新答案，就能很快定位错误。

\textbf{复杂度变化。} 暴力最坏 O(\ensuremath{n^{2}})，优化后每个 right 扫一次，哈希或数组访问均摊 O(1)，总时间 O(n)，空间取决于字符集。

\textbf{迁移追问。} 如果题目改成“若要求恰好包含 k 种字符，窗口条件从无重复变成种类计数”，先判断原状态是否仍然覆盖新答案集合，再决定是微调边界、改 value 语义，还是换算法。

\subsection{LeetCode 438 找到字符串中所有字母异位词}

\textbf{题目说明。} LeetCode 438 找到字符串中所有字母异位词 的题面入口要先还原成具体任务：固定长度窗口内字符频次等于目标频次即为答案。

\textbf{样例入口。} 拿 s=cbaebabacd、p=abc 手算，长度为 3 的窗口 cba 频次和 abc 相同，起点 0 是答案；窗口右移时只删除 c、加入 e。

\textbf{暴力基线。} 先统计 p 的频次 target，再枚举 s 中每个长度 k=p.size() 的窗口，为每个窗口重新统计 window 后比较。

\textbf{暴力过程。} 以 s=cbaebabacd、p=abc 为例，暴力窗口依次是 cba、bae、aeb、eba、bab、aba、bac、acd。每个窗口都重新数 k 个字符。

\textbf{重复工作。} 题面模型：固定长度窗口内字符频次等于目标频次即为答案。暴力版的慢点要落到本题：相邻窗口只差一个出去字符和一个进来字符。从 cba 到 bae，只有 c 出去、e 进来，b 和 a 不该重数。后面的优化只消掉这类重复工作，不能改变答案集合。

\textbf{优化条件。} 本题的关键观察是：右进左出维护计数，窗口长度达到目标长度后检查差异。这句话必须能翻译成可维护状态；如果题目条件改变后观察不再成立，就不能继续套原来的写法。

\textbf{相邻状态。} 规律是固定长度窗口维护字符频次差，重叠答案不能漏。把这个特例继续拆开看：右端进入只带来一个新元素，左端离开只删掉一个旧元素；只有当旧左端被移走后不可能再产生更优答案时，窗口才可以单向推进。若这个单调淘汰条件不存在，就要换成前缀和、队列或其他状态。

\textbf{状态复用。} k 来自 p.length()。先统计第一个窗口，之后每次窗口右移一格：window[s[i-1]]--，window[s[i+k-1]]++。手算时只改被新输入影响的那一小部分，而不是回头重算完整候选。

\textbf{指针和字段。} 状态字段要能说清：i 是当前窗口起点，out=s[i-1] 是离开字符，in=s[i+k-1] 是进入字符，target/window 是 26 长度频次数组。手算时按这个协议跟踪：第一个窗口 cba 与 abc 频次相同，记录 0；滑到 bae 后 c 计数减一、e 加一；滑到 bac 时频次再次相同，记录 6。

\textbf{代码落点。} 关键代码是固定窗口右进左出。若用 unordered\_map，计数减到 0 的 key 要 erase；若题目限定小写字母，用 array/vector 长度 26 更简单。关键代码行必须能逐句翻译成“旧状态怎样变成新状态”，否则就只是模板。

\textbf{正确性。} 每个被检查窗口长度都等于 k，window 频次始终精确表示当前窗口；频次等于 target 当且仅当该窗口是 p 的异位词。

\textbf{边界实验。} 先用目标比源串长、重复字符、字符集固定、答案重叠做反例检查。实验时重点跑目标比源串长、重复字符、字符集固定、答案重叠，并打印左右端、窗口计数和答案。若左端发生回退，或者窗口失去合法性后仍更新答案，就能很快定位错误。

\textbf{复杂度变化。} 暴力 O(nk)，固定窗口维护 O(n*26)，因字符集固定通常记作 O(n)，空间 O(26)。

\textbf{迁移追问。} 如果题目改成“维护差异计数可以避免每次比较整个数组”，先判断原状态是否仍然覆盖新答案集合，再决定是微调边界、改 value 语义，还是换算法。

\subsection{LeetCode 713 乘积小于 K 的子数组}

\textbf{题目说明。} LeetCode 713 乘积小于 K 的子数组 的题面入口要先还原成具体任务：正整数数组的窗口乘积随右端扩展变大、随左端收缩变小。

\textbf{样例入口。} 拿 nums=[10, 5, 2, 6]、k=100 手算，右端到 2 时乘积 100 不满足严格小于，左端移走 10 后窗口 [5, 2] 合法，以 2 结尾的合法子数组有 [2] 和 [5, 2] 两个。

\textbf{暴力基线。} 暴力枚举每个起点 left，向右连续乘积 product，product<k 就计数，否则停止或继续检查。

\textbf{暴力过程。} 以 [10, 5, 2, 6]、k=100 为例，从 10 开始可数 [10]、[10, 5]，到 [10, 5, 2] 乘积为 100 不满足；从 5 又重新乘 5、5*2、5*2*6。

\textbf{重复工作。} 题面模型：正整数数组的窗口乘积随右端扩展变大、随左端收缩变小。暴力版的慢点要落到本题：浪费在相邻起点重复计算乘积。数组是正整数，乘入右端会让乘积不减，除出左端会让乘积不增。后面的优化只消掉这类重复工作，不能改变答案集合。

\textbf{优化条件。} 本题的关键观察是：右端乘入新数，乘积不小于 k 时左端不断除出，合法后以 right 结尾的子数组数为 right-left+1。这句话必须能翻译成可维护状态；如果题目条件改变后观察不再成立，就不能继续套原来的写法。

\textbf{相邻状态。} 规律是正整数乘积窗口具有单调性，每个右端贡献 right-left+1 个合法起点。把这个特例继续拆开看：右端进入只带来一个新元素，左端离开只删掉一个旧元素；只有当旧左端被移走后不可能再产生更优答案时，窗口才可以单向推进。若这个单调淘汰条件不存在，就要换成前缀和、队列或其他状态。

\textbf{状态复用。} left/right 维护乘积 product。每次 product*=nums[right]；当 product>=k 时不断 product/=nums[left++]。合法后新增 right-left+1 个以 right 结尾的子数组。手算时只改被新输入影响的那一小部分，而不是回头重算完整候选。

\textbf{指针和字段。} 状态字段要能说清：product 是当前窗口乘积，left 是最小合法起点，right-left+1 是以 right 结尾的合法起点数量。手算时按这个协议跟踪：到 right=2 时 product=100，不满足，移走 10 得 product=10，窗口 [5, 2] 合法，新增 [2] 和 [5, 2] 两个。

\textbf{代码落点。} 关键是 k<=1 时直接返回 0，因为所有 nums 为正，乘积不可能严格小于 1。计数用 ans += right-left+1。关键代码行必须能逐句翻译成“旧状态怎样变成新状态”，否则就只是模板。

\textbf{正确性。} 在 product<k 时，当前窗口的任意后缀乘积也小于 k，所以以 right 结尾且起点在 [left, right] 的子数组全合法。

\textbf{边界实验。} 先用k 小于等于一、元素都为一、乘积溢出、窗口缩到空、严格小于而不是小于等于做反例检查。实验时重点跑k 小于等于一、元素都为一、乘积溢出、窗口缩到空、严格小于而不是小于等于，并打印左右端、窗口计数和答案。若左端发生回退，或者窗口失去合法性后仍更新答案，就能很快定位错误。

\textbf{复杂度变化。} 暴力 O(\ensuremath{n^{2}})，滑动窗口 O(n) 时间、O(1) 空间。

\textbf{迁移追问。} 如果题目改成“若数组含零或负数，乘积单调性被破坏，需要换模型”，先判断原状态是否仍然覆盖新答案集合，再决定是微调边界、改 value 语义，还是换算法。

\subsection{LeetCode 862 和至少为 K 的最短子数组}

\textbf{题目说明。} LeetCode 862 和至少为 K 的最短子数组 的题面入口要先还原成具体任务：数组允许负数，普通滑动窗口的和不再随左端移动单调变化。

\textbf{样例入口。} 拿 [2, -1, 2]、K=3 手算，普通正数窗口会被 -1 破坏单调性；前缀和为 0, 2, 1, 3，只有 3-0 达到 3。

\textbf{暴力基线。} 暴力枚举所有 left/right，用前缀和 O(1) 求区间和，检查是否至少为 K 并更新最短长度。

\textbf{暴力过程。} 以 [2, -1, 2]、K=3 为例，前缀和是 [0, 2, 1, 3]。普通正数窗口会被 -1 破坏：右端增加后和可能变小。

\textbf{重复工作。} 题面模型：数组允许负数，普通滑动窗口的和不再随左端移动单调变化。暴力版的慢点要落到本题：浪费在每个右端都回头扫描所有历史左端。需要的是某些前缀和尽量小且下标尽量靠后的候选左端。后面的优化只消掉这类重复工作，不能改变答案集合。

\textbf{优化条件。} 本题的关键观察是：在前缀和数组上用单调队列维护可能成为最优左端的下标。这句话必须能翻译成可维护状态；如果题目条件改变后观察不再成立，就不能继续套原来的写法。

\textbf{相邻状态。} 规律是在前缀和上用单调队列维护可能的最优左端，队首能满足时结算长度。把这个特例继续拆开看：右端进入只带来一个新元素，左端离开只删掉一个旧元素；只有当旧左端被移走后不可能再产生更优答案时，窗口才可以单向推进。若这个单调淘汰条件不存在，就要换成前缀和、队列或其他状态。

\textbf{状态复用。} 构造前缀和 prefix。用单调递增队列保存可能作为左端的前缀下标；当前前缀减队首前缀 >=K 时更新答案并弹出队首；队尾前缀不小于当前前缀时被当前支配，弹出。手算时只改被新输入影响的那一小部分，而不是回头重算完整候选。

\textbf{指针和字段。} 状态字段要能说清：deque 中保存前缀下标，且对应 prefix 值递增；i 是当前右端后一位，i-deque.front 是候选长度。手算时按这个协议跟踪：prefix 到 3 时，3-0>=3，长度 3；若出现更小的新前缀，它会支配队尾较大前缀，因为未来用它能得到更大区间和且更短或不更差。

\textbf{代码落点。} 关键有两个 while：先从队首结算满足 K 的最短长度，再从队尾删除 prefix 值不小于当前值的被支配候选。关键代码行必须能逐句翻译成“旧状态怎样变成新状态”，否则就只是模板。

\textbf{正确性。} 队首是当前最早且前缀较小的候选；一旦满足 K，继续保留它只会让未来长度更长，所以可以弹出。队尾被更小且更新的前缀支配，永远不会更优。

\textbf{边界实验。} 先用负数、答案不存在、K 为负或零、队尾支配关系做反例检查。实验时重点跑负数、答案不存在、K 为负或零、队尾支配关系，并打印左右端、窗口计数和答案。若左端发生回退，或者窗口失去合法性后仍更新答案，就能很快定位错误。

\textbf{复杂度变化。} 暴力 O(\ensuremath{n^{2}})，单调队列中每个前缀入队出队各一次，时间 O(n)，空间 O(n)。

\textbf{迁移追问。} 如果题目改成“这是从窗口题升级到前缀和加单调队列的代表”，先判断原状态是否仍然覆盖新答案集合，再决定是微调边界、改 value 语义，还是换算法。

\begin{principlebox}{本节复盘方式}
不要只看最终算法名。每道题复盘时先遮住优化版，口述暴力枚举对象；再指出相邻窗口、历史前缀、候选区间、搜索状态或子问题里哪一部分被重复处理；最后用一个边界样例验证状态更新顺序。能讲完这三步，才算真正掌握本章案例。
\end{principlebox}
